\chap{智能交互与工作投影：让理解和控制共处}
\label{sec:interaction}
\subsection{信息集中之后仍然存在理解成本}
WALL 的实践说明，集中保存设计信息只解决了可查找性的一部分问题。面对几千个设计点，人需要判断哪些值得看；面对一句缺少背景的说明，人又需要重新建立语境。AI 在生成这句话时可能掌握实现、历史和依赖，而读者只看见结论。压缩文本减少了字数，却可能增加解释成本。

实际软件界面提供了另一种条件。人正在尝试某项任务，能够看见对象、操作和结果，因而可以直接指出哪里不符合预期。把这种能力扩展到开发过程，意味着为架构、协议、方案和进展提供同样具体的理解入口。视图的价值在于支持判断和行动，而不只在于展示更多数据。

因此，智能交互系统需要回答：当前人想理解什么，已经知道什么，还缺少哪些关键关系，以及理解后能够怎样参与。表达形式随问题改变。架构可能适合层次图，行为适合可操作演示，方案差异适合共同情境下的对比，进展则适合突出自上次查看后的重要变化。

\subsection{工作语义与交互投影}
工作语义包括目标、职责、模块、协议、任务、证据和结果等对象及其关系。交互投影根据人的当前问题选择其中一部分，组织成可理解的视图和操作。Lineage 可以承载团队与工作关系，Workbench 将相关内容投影为人的工作界面。软件架构、组织关系和运行状态仍有各自含义，需要通过明确关联组合。

投影过程可以分为信息选择、解释组织、媒介选择、操作绑定和结果更新。信息选择确定当前重点，解释组织提供背景与因果关系，媒介选择决定图、文字、演示或其他形式，操作绑定把人的动作关联到真实工作对象，结果更新则在工作变化后刷新相应表达。

\fig{09-projection}{交互投影的生成与回传。视图来自工作对象和人的关注点，人的操作携带对象与参数进入实际执行，界面随后依据真实结果更新。}{projection}

同一个对象可以在多个视图中出现。任务取消协议可以在架构图中表现为模块之间的边，在机制视图中表现为事件序列，在用户旅程中表现为状态变化。对象身份和关键语义应保持一致，使人能够沿着疑问切换视图而不丢失背景。各视图不必显示相同细节，但应避免对实际状态给出相互矛盾的解释。

\subsection{理解与控制发生在同一情境中}
当人比较两个方案时，界面可以允许切换相同场景，查看行为差异，然后在该对象上表达偏好。发现某一步不符合预期时，可以标注那一步并说明期望。反馈自动携带场景、对象、当前状态和可见结果，减少人重新组织长段描述的负担。

这种设计也能改善需求理解。例如，人体验取消演示后说“我希望看到停止确认之前的等待状态”，系统能够把反馈联系到当前状态映射，而不必猜测“那里”指什么。若反馈涉及软件行为，转入相应开发任务；若仅涉及解释形式，更新投影规则。两类工作可能相关，但需要知道实际改变了什么。

混合主动性交互研究讨论自动化服务与直接操作的结合\cite{horvitz1999}。在本文的场景中，Agent 可以主动选择表达和提出建议，人则能够在需要时进入对象、调整约束和改变方向。自主性减少低价值的微观指挥，可操作表达帮助人保留关键控制。

\subsection{表达必须说明状态与作用}
生成界面可能看起来完整，却展示尚未实施的方案。因此需要让当前实现、拟议设计、模拟行为和已验证结果在相关位置清楚可辨。一个模拟取消按钮应说明它作用于演示环境；真实任务上的取消操作应关联实际任务，并依据真实执行状态更新。视觉相似不能代替语义准确。

状态说明也应避免过度承诺。请求已发送、服务已接收、执行器已停止和资源已释放，可能是不同事实。对用户有意义的界面可以简化细节，但保留下来的解释应足以支持正确判断。若某项结果尚未确认，可以呈现等待及可用的后续操作。

重要结论需要提供按需查看依据的入口。默认界面不必充满日志，但人应能在疑问出现时追查相关检查和实际对象。这种渐进式呈现与信息可视化中由全貌进入细节的思路相呼应\cite{shneiderman1996}。关键在于选择哪些细节有助于当前任务，而非简单折叠所有信息。

\subsection{投影如何保持准确与连续}
投影生成至少需要核对事实、语义和操作三个方面。事实核对检查当前状态是否来自有效记录；语义核对检查图示和说明是否保留了关键区别；操作核对检查按钮实际作用于哪个对象，以及结果怎样返回界面。三者中任一部分脱节，都可能让人形成错误判断。

例如，架构图可能仍显示旧的模块归属，但运行协议已经迁移。界面可以在相应对象上显示更新时间和需要刷新之处，并从实际对象获取新的关系。若人正在比较一个尚未实施的方案，则保留该方案的固定背景，明确它与当前运行状态的差异。持续更新并不要求把所有视图无条件切换成最新状态。

为了降低学习成本，可以保持对象名称、导航层次和基本动作的稳定性，将变化集中在当前问题的表达上。一个人多次进入任务服务后，应能找到相同职责与协议入口；新的演示则解释近期变化。这样，界面的适应性与长期方向感能够共同发展。

\subsection{多媒体选择与生成成本}
网页能够组合交互控件、状态视图和直接反馈；PPT 适合组织阶段成果与讲解顺序；视频适合展示时间变化和操作过程；3D 适合需要空间理解的对象。媒体选择应围绕理解任务，不因技术可用就增加展示复杂度。某个协议用几步可暂停动画就能解释清楚时，复杂三维场景未必带来收益。

多媒体还需要提供可访问的替代表达。动态图应允许暂停和回看，关键内容配有文字说明，颜色之外保留标签与结构差异。跨设备查看时，字号、布局和操作区域需要重新考虑。生成界面的质量包含这些使用条件，而不只是单张截图的美观。

生成成本可以通过复用稳定组件、局部更新和按需展开控制。工作事实变化时，受影响的视图需要更新，但不必重新制作全部展示。对很少被查看的细节，可以在用户进入时生成；对关键状态，则应保持及时可见。成本分配由使用频率、重要性和生成延迟共同影响。

\subsection{异步汇报与表达的持续演化}
人不会持续在线，汇报本身也需要时间。系统可以委派独立成员在阶段工作之后整理成果、更新界面和准备解释。人再次进入时，首先获得与上次查看相比的重要变化，再决定深入哪里。为了避免信息堆积，系统需要区分已经理解的内容、仍需回应的问题和新出现的影响。

界面可以通过反馈持续改进。人说“协议以后用具体交互来解释”，可以形成该类任务的表达偏好；人指出某个动态图过快，则调整播放与逐步观察方式。偏好应保存适用范围，允许后续纠正。不同读者和任务可能需要不同表达，个性化不应让核心对象失去共同含义。

人机 AI 交互指南强调能力说明、纠错与随时间变化的互动\cite{amershi2019}。本文据此提出一个工程化方向：表达生成作为持续工作，接收与软件开发相同的实际反馈。开发系统既改善软件，也改善人参与软件发展的方式。

\begin{insight}{投影的成功标准}
人能够更快理解当前情况，在合适对象上表达准确反馈，并看见操作对实际工作的作用。界面的内容数量、动画数量和生成 token 数，都需要通过这些结果解释其价值。
\end{insight}
